% Chapter 6: Infinite Products and the Sine Function
% ch06_infinite_products.tex

\chapter{Infinite Products: Wallis, Sinc, Vi\`{e}te, and the Sine Function}

\begin{historical}
Euler's work on infinite products was revolutionary. He discovered the product formula for
the sine function, the Wallis product for $\pi/2$, Vi\`{e}te's formula for $\pi$, and the Euler
product for the zeta function. These products established deep connections between analysis,
number theory, and geometry that continue to drive research today.
\end{historical}

\section{The Wallis Product}

John Wallis discovered (1655) and Euler refined the product formula for $\pi/2$:
\[
\frac{\pi}{2} = \prod_{n=1}^\infty \frac{4n^2}{4n^2 - 1}
= \frac{2}{1} \cdot \frac{2}{3} \cdot \frac{4}{3} \cdot \frac{4}{5} \cdot \frac{6}{5} \cdot \frac{6}{7} \cdots
\]

\begin{theorem}[Wallis Product]
\[
\frac{\pi}{2} = \prod_{n=1}^{\infty} \frac{2n}{2n-1} \cdot \frac{2n}{2n+1}
\]
\end{theorem}

\begin{proof}
Starting from the beta function $B(m,n) = \int_0^1 x^{m-1}(1-x)^{n-1} dx$, one can show
that $\int_0^{\pi/2} \sin^{2n} x \, dx = \frac{(2n)!}{2^{2n}(n!)^2} \cdot \frac{\pi}{2}$ and
$\int_0^{\pi/2} \sin^{2n+1} x \, dx = \frac{2^{2n}(n!)^2}{(2n+1)!}$. Taking the ratio and
letting $n \to \infty$ gives the Wallis product.
\end{proof}

\subsection{Proof via Gamma Function}

\begin{theorem}
Using $\Gamma(n+1) = n!$ and $\Gamma(1/2) = \sqrt{\pi}$:
\[
\frac{\pi}{2} = \lim_{n\to\infty} \frac{2^{4n} (n!)^4}{[(2n)!]^2 (2n+1)}
\]
\end{theorem}

\begin{proof}
From the duplication formula $\Gamma(z)\Gamma(z+1/2) = 2^{1-2z}\sqrt{\pi}\Gamma(2z)$:
\[
\Gamma(n+1)\Gamma(n+1/2) = 2^{-2n}\sqrt{\pi}\Gamma(2n+1)
\]
Setting $z = n+1/2$ and taking ratios yields the Wallis product after simplification.
\end{proof}

\begin{lstlisting}[caption={Wallis product for $\pi/2$}, label=lst:wallis]
from euler import wallis_product, wallis_product_partial
import math

print("$\pi$ approximation via Wallis product:")
for n in [1, 2, 5, 10, 50, 100, 1000]:
    approx = wallis_product(n) * 2
    print(f"  n={n:4d}: $\pi$ \approx {approx:.12f}  error = {abs(math.pi - approx):.2e}")
\end{lstlisting}

\section{Convergence Analysis of Infinite Products}

\subsection{Necessary and Sufficient Conditions}

\begin{theorem}[Convergence Criterion]
The infinite product $\prod_{n=1}^{\infty}(1 + a_n)$ converges (to a nonzero value) if and only if
$\sum_{n=1}^{\infty} \ln(1 + a_n)$ converges.
\end{theorem}

\begin{theorem}[Absolute Convergence]
The product $\prod(1 + a_n)$ converges absolutely if and only if $\sum |a_n| < \infty$.
\end{theorem}

\subsection{Comparison Tests}

\begin{theorem}
If $a_n \geq 0$, then $\prod(1 + a_n)$ converges if and only if $\sum a_n$ converges.
\end{theorem}

\begin{theorem}
If $a_n \to 0$, then $\prod(1 + a_n)$ converges if and only if $\sum a_n$ converges.
\end{theorem}

\subsection{Convergence Rate Classification}

\begin{table}[H]
\centering
\begin{tabular}{l|c|c}
Product Type & Convergence Rate & Terms for $\epsilon$ Accuracy \\
\hline
Geometric ($|r|<1$) & $O(r^n)$ & $O(\log(1/\epsilon))$ \\
Power law ($a_n = O(n^{-\alpha})$) & $O(n^{1-\alpha})$ & $O(\epsilon^{-1/(1-\alpha)})$ \\
Double exponential & Super-fast & $O(\log\log(1/\epsilon))$ \\
\hline
\end{tabular}
\caption{Convergence rates for common infinite products}
\end{table}

\section{The Sinc Product}

Euler's product for $\sin(\pi x)/(\pi x)$ is a fundamental tool:
\[
\frac{\sin(\pi x)}{\pi x} = \prod_{n=1}^\infty \left(1 - \frac{x^2}{n^2}\right)
\]

This formula follows from the Weierstrass factorization theorem: an entire function is
determined (up to an exponential factor) by its zeros. Since $\sin(\pi x)$ has zeros at
all integers, the infinite product over $(1 - x/n)$ captures all zeros.

\begin{theorem}[Euler's Sinc Product]
For all complex $x$:
\[
\frac{\sin(\pi x)}{\pi x} = \prod_{n=1}^{\infty} \left(1 - \frac{x^2}{n^2}\right)
\]
\end{theorem}

\begin{proof}
Consider $f(x) = \frac{\sin(\pi x)}{\pi x}$. This is an entire function with simple zeros
at all nonzero integers $x = \pm n$ for $n \geq 1$. By the Weierstrass factorization theorem
for entire functions of order 1, $f$ can be written as:
\[
f(x) = e^{g(x)} \prod_{n=1}^{\infty} \left(1 - \frac{x}{n}\right)\left(1 + \frac{x}{n}\right)
= e^{g(x)} \prod_{n=1}^{\infty} \left(1 - \frac{x^2}{n^2}\right)
\]
Since $f(0) = 1$ and $g(x)$ must be constant for convergence, we get $g(x) = 0$, yielding
the result.
\end{proof}

\subsection{Weierstrass Factorization Theorem}

\begin{theorem}[Weierstrass]
Every entire function can be factored into a product involving its zeros:
\[
f(z) = z^m e^{g(z)} \prod_{n=1}^{\infty} E_{p_n}\left(\frac{z}{a_n}\right)
\]
where $E_k(z) = (1-z)\exp(z + z^2/2 + \cdots + z^k/k)$ are the Weierstrass elementary factors.
\end{theorem}

\begin{lstlisting}[caption={Sinc product}, label=lst:sinc]
from euler import sinc_product, sinc_product_partial
import math

for x in [0.1, 0.25, 0.5]:
    exact = math.sin(math.pi * x) / (math.pi * x)
    print(f"\n  x = {x}: exact = {exact:.10f}")
    for n in [10, 50, 100, 500]:
        approx = sinc_product(x, n)
        print(f"    n={n:3d}: {approx:.10f}  error = {abs(exact - approx):.2e}")
\end{lstlisting}

\section{Euler's Product for Sine}

Setting $x \to x/\pi$ in the sinc product gives:
\[
\sin x = x \prod_{n=1}^\infty \left(1 - \frac{x^2}{n^2 \pi^2}\right)
\]

This was Euler's key insight for solving the Basel problem.

\subsection{Convergence Analysis}

\begin{theorem}
The product $\prod_{n=1}^{\infty}\left(1 - \frac{x^2}{n^2\pi^2}\right)$ converges uniformly on compact subsets of $\mathbb{C}$.
\end{theorem}

\begin{proof}
For fixed $x$ and large $n$, we have $|x^2/(n^2\pi^2)| < 1/2$. Using the estimate
$|\ln(1 - a_n)| \leq 2|a_n|$ for small $a_n$:
\[
\sum_{n=1}^{\infty} \left|\frac{x^2}{n^2\pi^2}\right| = \frac{|x|^2}{\pi^2}\sum_{n=1}^{\infty}\frac{1}{n^2} = \frac{|x|^2}{6} < \infty
\]
By the Weierstrass M-test, the product converges absolutely and uniformly.
\end{proof}

\begin{lstlisting}[caption={Euler's sine product}, label=lst:euler_sine]
from euler import euler_product_sine
import math

for x in [0.5, 1.0, 2.0, math.pi/2]:
    exact = math.sin(x)
    print(f"\n  x = {x:.4f}: sin(x) = {exact:.8f}")
    for n in [10, 50, 100]:
        approx = euler_product_sine(x, n)
        print(f"    n={n:3d}: {approx:.8f}  error = {abs(exact - approx):.2e}")
\end{lstlisting}

\section{Product for Cosine}

Euler also discovered the product formula for cosine by shifting the sine product:
\[
\cos x = \prod_{n=1}^\infty \left(1 - \frac{4x^2}{(2n-1)^2 \pi^2}\right)
\]

This follows from $\cos x = \sin(x + \pi/2)$ and the sine product formula.

\subsection{Proof of Cosine Product}

\begin{theorem}
\[
\cos x = \prod_{n=1}^{\infty} \left(1 - \frac{4x^2}{(2n-1)^2\pi^2}\right)
\]
\end{theorem}

\begin{proof}
Using $\cos x = \sin(x + \pi/2)$ and the sine product:
\[
\cos x = \left(x + \frac{\pi}{2}\right) \prod_{n=1}^{\infty} \left(1 - \frac{(x+\pi/2)^2}{n^2\pi^2}\right)
\]
Substituting $n \to n - 1/2$ and simplifying yields the product over odd integers.
\end{proof}

\section{Vi\`{e}te's Formula for $\pi$}

Vi\`{e}te's formula (1593) was the first infinite product in mathematical history:
\[
\frac{2}{\pi} = \prod_{n=1}^\infty \cos\left(\frac{\pi}{2^{n+1}}\right)
= \cos\frac{\pi}{4} \cdot \cos\frac{\pi}{8} \cdot \cos\frac{\pi}{16} \cdots
\]

\begin{theorem}[Vi\`{e}te's Formula]
\[
\frac{2}{\pi} = \prod_{n=1}^{\infty} \cos\left(\frac{\pi}{2^{n+1}}\right)
\]
\end{theorem}

\begin{proof}
Starting from the double-angle formula $\sin(2x) = 2\sin x \cos x$, we get
$\cos x = \frac{\sin(2x)}{2\sin x}$. Setting $x = \pi/4$ and iterating:
\[
\cos\frac{\pi}{4} \cdot \cos\frac{\pi}{8} \cdot \cos\frac{\pi}{16} \cdots
= \frac{\sin(\pi/2)}{2} \cdot \frac{\sin(\pi/4)}{2\sin(\pi/8)} \cdot \frac{\sin(\pi/4)}{2\sin(\pi/8)} \cdots
\]
The telescoping product yields $2/\pi$.
\end{proof}

\subsection{Nested Radical Form}

Vi\`{e}te's formula can be written using nested radicals:
\[
\frac{2}{\pi} = \frac{\sqrt{2}}{2} \cdot \frac{\sqrt{2+\sqrt{2}}}{2} \cdot \frac{\sqrt{2+\sqrt{2+\sqrt{2}}}}{2} \cdots
\]

\begin{lstlisting}[caption={Vi\`{e}te's formula for $\pi$}, label=lst:viete]
from euler import viete_formula
import math

print("Vi\`{e}te's formula:")
for n in [5, 10, 20, 30, 50]:
    approx = viete_formula(n)
    print(f"  n={n:2d}: $\pi$ \approx {approx:.12f}  error = {abs(math.pi - approx):.2e}")
\end{lstlisting}

\section{Euler Product for the Zeta Function}

For $\Re(s) > 1$:
\[
\zeta(s) = \prod_{p} \frac{1}{1 - p^{-s}}
\]

where the product is over all primes $p$.

\subsection{Proof of Euler Product}

\begin{theorem}[Euler Product Formula]
For $\Re(s) > 1$:
\[
\zeta(s) = \sum_{n=1}^{\infty} \frac{1}{n^s} = \prod_{p \text{ prime}} \left(1 - \frac{1}{p^s}\right)^{-1}
\]
\end{theorem}

\begin{proof}
Using the geometric series $\frac{1}{1 - p^{-s}} = \sum_{k=0}^{\infty} p^{-ks}$:
\[
\prod_{p} \left(1 - \frac{1}{p^s}\right)^{-1} = \prod_{p} \sum_{k=0}^{\infty} \frac{1}{p^{ks}}
\]
By the fundamental theorem of arithmetic, every integer $n$ has a unique prime factorization,
so expanding the product gives exactly $\sum_{n=1}^{\infty} n^{-s} = \zeta(s)$.
\end{proof}

\subsection{Consequences}

\begin{theorem}
The Euler product implies there are infinitely many primes.
\end{theorem}

\begin{proof}
If there were finitely many primes, the product would be finite, but $\zeta(s) \to \infty$
as $s \to 1^+$, contradicting finiteness.
\end{proof}

\begin{lstlisting}[caption={Euler product for $\zeta(2)$}, label=lst:euler_zeta_product]
from euler import euler_product_formula
import math

print("Euler product for $\zeta(2) = \pi^2/6$:")
exact = math.pi**2 / 6
for n in [10, 50, 100, 500, 1000]:
    approx = euler_product_formula(2, n)
    print(f"  {n:4d} primes: {approx:.10f} (error: {abs(exact-approx):.2e})")
\end{lstlisting}

\section{Computational Analysis}

\subsection{Convergence Rates}

\begin{table}[H]
\centering
\begin{tabular}{l|c|c}
Product & Convergence Rate & Approximate terms for 6 digits \\
\hline
Wallis ($\pi/2$) & $O(1/n)$ & $\sim 10^6$ \\
Vi\`{e}te ($\pi$) & Doubles digits per term & $\sim 20$ \\
Sinc product & $O(1/n)$ & $\sim 10^6$ \\
Euler zeta product & $O(1/p_N)$ & $\sim 10^5$ \\
\hline
\end{tabular}
\caption{Convergence rates for famous infinite products}
\end{table}

\subsection{Error Bounds}

\begin{theorem}
For the Wallis product after $n$ terms:
\[
\left|\frac{\pi}{2} - \prod_{k=1}^{n} \frac{4k^2}{4k^2-1}\right| \sim \frac{1}{4n}
\]
\end{theorem}

\begin{theorem}
For Vi\`{e}te's formula after $n$ terms:
\[
\left|\pi - \pi_n\right| = O(4^{-2^n})
\]
\end{theorem}

\subsection{Complexity Analysis}

\begin{itemize}
\item \textbf{Wallis product:} Each term requires 2 multiplications and 1 division. Total: $O(n)$ operations for $n$ terms.
\item \textbf{Vi\`{e}te formula:} Each term requires one cosine evaluation. With nested radicals, the arithmetic complexity is $O(n^2)$ bit operations for $n$ terms.
\item \textbf{Sinc product:} Each term requires 2 multiplications. Total: $O(n)$ operations.
\item \textbf{Euler zeta product:} Computing $n$ primes takes $O(n \log n \log\log n)$ with sieve. Each term requires one exponentiation.
\end{itemize}

\section{Deriving the Basel Problem}

Setting $x = 1/2$ in the sinc product:
\[
\frac{\sin(\pi/2)}{\pi/2} = \frac{2}{\pi} = \prod_{n=1}^\infty \left(1 - \frac{1}{4n^2}\right)
\]
Inverting gives the Wallis product. Expanding the sinc product as a power series and
comparing coefficients yields $\zeta(2) = \pi^2/6$.

\subsection{Taylor Expansion of Sinc}

\begin{theorem}
Expanding $\frac{\sin(\pi x)}{\pi x}$ as a power series:
\[
\frac{\sin(\pi x)}{\pi x} = 1 - \frac{\pi^2 x^2}{3!} + \frac{\pi^4 x^4}{5!} - \cdots
\]
\end{theorem}

\begin{proof}
Using the Taylor series $\sin(\pi x) = \pi x - \frac{(\pi x)^3}{3!} + \frac{(\pi x)^5}{5!} - \cdots$:
\[
\frac{\sin(\pi x)}{\pi x} = 1 - \frac{(\pi x)^2}{3!} + \frac{(\pi x)^4}{5!} - \cdots
\]
\end{proof}

\subsection{Matching Coefficients}

Expanding the product:
\[
\prod_{n=1}^{\infty}\left(1 - \frac{x^2}{n^2}\right) = 1 - \left(\sum_{n=1}^{\infty}\frac{1}{n^2}\right)x^2 + \cdots
\]

Comparing coefficients of $x^2$:
\[
-\frac{\pi^2}{6} = -\sum_{n=1}^{\infty}\frac{1}{n^2} = -\zeta(2)
\]

\section{Riemann Xi Function}

\begin{definition}
The Riemann Xi function is defined as:
\[
\xi(s) = \frac{1}{2}s(s-1)\pi^{-s/2}\Gamma\left(\frac{s}{2}\right)\zeta(s)
\]
\end{definition}

\begin{theorem}
$\xi(s)$ satisfies the functional equation $\xi(s) = \xi(1-s)$ and has an Euler product:
\[
\xi(s) = \xi(0) \prod_{\rho} \left(1 - \frac{s}{\rho}\right)
\]
where $\rho$ runs over the non-trivial zeros of $\zeta(s)$.
\end{theorem}

\section{Frequently Asked Questions}

\begin{itemize}
\item \textbf{Q: Why does the Wallis product converge so slowly?} \\
The Wallis product converges like $1/n$. After 1,000 terms, the approximation of $\pi/2$ has only about 2--3 correct decimal places. This is why it is historically significant but practically useless for computing $\pi$.

\item \textbf{Q: How does the sinc product relate to the Basel problem?} \\
Setting $x$ in the sinc product $\frac{\sin(\pi x)}{\pi x} = \prod_{n=1}^{\infty} (1 - x^2/n^2)$ and expanding both sides gives coefficients involving $\sum 1/n^k$. Comparing the $x^2$ coefficient yields $\zeta(2) = \pi^2/6$, the Basel result. The $x^4$ coefficient gives $\zeta(4) = \pi^4/90$.

\item \textbf{Q: Why does the sine product vanish at $x = n\pi$?} \\
At $x = n\pi$ for any nonzero integer $n$, the factor $(1 - n^2\pi^2/(n^2\pi^2)) = 0$ appears in the product, making the entire product zero. This matches $\sin(n\pi) = 0$.

\item \textbf{Q: What is the difference between the Wallis product and Vi\`{e}te's formula?} \\
The Wallis product $\pi/2 = \prod_{n=1}^{\infty} \frac{4n^2}{4n^2-1}$ involves rational factors and converges slowly. Vi\`{e}te's formula $\frac{2}{\pi} = \prod_{n=1}^{\infty} \cos(\frac{\pi}{2^{n+1}})$ involves nested radicals and converges much faster (doubling the number of correct digits with each factor).

\item \textbf{Q: Can infinite products represent other elementary functions?} \\
Yes. $\cos x$, $\frac{\sin x}{x}$, the gamma function, and the Riemann zeta function all have product representations. Euler discovered many of these, and they remain fundamental tools in analytic number theory and complex analysis.

\item \textbf{Q: Why does $\prod_{n=1}^{\infty} (1 - 1/n^2) = 1/2$?} \\
The partial product telescopes: $\prod_{n=2}^N (1 - 1/n^2) = \prod_{n=2}^N \frac{(n-1)(n+1)}{n^2} = \frac{1 \cdot (N+1)}{2N} \to \frac{1}{2}$ as $N \to \infty$. The $n=1$ term is excluded because $1 - 1/1^2 = 0$.

\item \textbf{Q: What is the connection between infinite products and prime numbers?} \\
The Euler product $\zeta(s) = \prod_p (1 - p^{-s})^{-1}$ directly connects the Riemann zeta function to the distribution of prime numbers. This connection is the foundation of analytic number theory.
\end{itemize}

\section{Exercises}

\begin{exercise}
Prove the Wallis product by integrating $\int_0^{\pi/2} \sin^n x \, dx$ and taking the ratio of consecutive integrals.
\end{exercise}

\begin{exercise}
Derive the product formula for $\cos x$ from the sine product by using $\cos x = \sin(x + \pi/2)$.
\end{exercise}

\begin{exercise}
Show that $\prod_{n=1}^{\infty} (1 - 1/n^2) = 1/2$ by telescoping.
\end{exercise}

\begin{exercise}[Proof]
Use the sinc product to derive $\zeta(2) = \pi^2/6$ by comparing coefficients in the power series expansion.
\end{exercise}

\begin{exercise}[Computational]
Implement Vi\`{e}te's formula using nested radicals and verify it converges faster than the Wallis product. Count how many terms are needed to reach 6 decimal places of accuracy.
\end{exercise}

\begin{exercise}
Prove Vi\`{e}te's formula $2/\pi = \prod_{k=1}^{\infty} \cos(\pi/2^{k+1})$ using the double-angle formula $\sin(2x) = 2\sin x \cos x$.
\end{exercise}

\begin{exercise}
Show that the Euler product for $\zeta(s)$ at $s=2$ converges to $\pi^2/6$ using the first 100 primes. Compare the error with the series sum $\sum_{n=1}^{1000} 1/n^2$.
\end{exercise}

\begin{exercise}
Prove that $\prod_{n=1}^{\infty} \cos(x/2^n) = \frac{\sin x}{x}$ using the double-angle formula repeatedly.
\end{exercise}

\begin{exercise}
Show that the infinite product $\prod_{n=1}^{\infty} \left(1 + \frac{1}{n(n+2)}\right)$ converges to $3/2$.
\end{exercise}

\begin{exercise}[Proof]
Prove the Weierstrass factorization theorem for the sine function: $\sin(\pi z) = \pi z \prod_{n=1}^{\infty} (1 - z^2/n^2)$.
\end{exercise}
